\documentclass[11pt]{article}
\usepackage[margin=1in]{geometry}
\usepackage{amsmath,amssymb,amsfonts}
\usepackage{booktabs}
\usepackage{hyperref}

\title{NBGA: No-Backprop Gradient Approximation\\Training Residual Networks Without the Chain Rule}
\author{Taran S. Marley\\[4pt]\small\url{https://github.com/Benzidrine/ngba_training}}
\date{}

\begin{document}
\maketitle

\begin{abstract}
We propose No-Backprop Gradient Approximation (NBGA), a method to train residual neural networks where every layer receives the same final-output gradient $\delta_L$ as its learning signal. The chain rule is not computed between layers. The approximation rests on a simple observation: when residual weights are small, the Jacobian of the skip connection dominates the Jacobian of the residual function, making the gradient approximately constant across all layers. We demonstrate NBGA by fine-tuning Qwen3.5-0.8B (24 blocks, 752M parameters) on WikiText-103 and Alpaca, showing a 23\% perplexity improvement and successful instruction-following behaviour---all achieved with zero gradient passing between layers. We further scale NBGA to Qwen3.5-4B (32 blocks, 4B parameters), achieving a 32\% perplexity reduction and a +5.5\% improvement on HumanEval pass@1 using only $\delta_L$ broadcast with no proximal term. A proximal extension enables stable fine-tuning of pre-trained models. The method is parallelizable, memory-efficient, and scales with depth without gradient vanishing.
\end{abstract}

\section{Background}

Standard backpropagation computes the gradient of the loss $L$ with respect to each layer's parameters via the chain rule. For a residual block

\[
h_{l} = h_{l-1} + f_l(h_{l-1}),
\]

the gradient of the loss with respect to the block's input is

\[
\delta_{l-1} = \frac{\partial L}{\partial h_{l-1}} = \delta_l \cdot (I + J_{f_l})^T,
\]

where $J_{f_l} = \partial f_l / \partial h_{l-1}$ is the Jacobian of the residual function. This computation is sequential: $\delta_{l-1}$ requires $\delta_l$, which requires $\delta_{l+1}$, and so on. For $D$ layers, backprop performs $D$ sequential matrix multiplications that cannot be parallelized. In deep networks, repeatedly multiplying by $(I + J_f)^T$ causes gradients to vanish or explode exponentially with $D$.

\section{The NBGA Approximation}

When the residual function $f_l$ has small weights, $\|f_l(h)\| \ll \|h\|$, making $J_{f_l} \approx 0$. Consequently $(I + J_{f_l})^T \approx I$, and the gradient propagation simplifies to

\[
\delta_{l-1} \approx \delta_l \approx \delta_{l+1} \approx \cdots \approx \delta_L.
\]

The gradient at \emph{every} layer is approximately equal to the gradient at the final layer $\delta_L = \partial L / \partial h_L$. The weight gradient becomes a parallelizable outer product:

\[
\frac{\partial L}{\partial W_l} \approx h_{l-1} \otimes \delta_L.
\]

This involves no chain rule, no sequential computation, and no knowledge of any other layer's state. Each layer's gradient can be computed independently and in parallel.

\subsection{The Proximal Extension}

For pre-trained models, the weights are large and $\|J_f\|$ may not be small enough for $\delta_L \approx \delta_{\text{true}}$ to hold accurately. We introduce a proximal term $\lambda\|W - W_0\|^2$ that pulls each block's weights toward their pre-trained origin:

\[
W_l \gets W_l - \eta \cdot \nabla_{\text{NBGA}} - \lambda \cdot (W_l - W_l^{(0)}).
\]

The proximal term bounds cumulative approximation error. We use $\lambda = 10^{-5}$ throughout our experiments.

\section{Algorithm}

\begin{enumerate}
\item \textbf{Full forward:} Compute all block outputs, caching intermediate activations (or using reversible reconstruction).
\item \textbf{Final loss:} Compute $\delta_L = \partial L / \partial h_L$ from the head gradient.
\item \textbf{Per-block backward:} For each block, re-forward with detached input, backward with $\delta_L$:
   \[\frac{\partial L}{\partial W_l} = h_{l-1} \otimes \delta_L.\]
\item \textbf{Update:} Apply gradients with proximal term.
\end{enumerate}

The backward passes are independent across blocks and can be computed in any order.

\section{Results: Qwen3.5-0.8B Fine-tuning}

We further fine-tune Qwen3.5-0.8B (the instruct-tuned variant; 24 blocks, 752M parameters, 1024 hidden dimension) using Proximal NBGA on a single NVIDIA RTX 4090 (24GB). The base model is already capable of instruction-following; our NBGA fine-tuning adds additional data on top.

\begin{table}[h]
\centering
\begin{tabular}{lcccc}
\toprule
Stage & Dataset & Steps & Val Loss & $\Delta$ \\
\midrule
WikiText adaptation & WikiText-103 (5K) & 500 & --- & +1377 PPL (23\%) \\
Instruction tuning & Alpaca (2K) & 500 & 0.42 & 24.3\% \\
GPT-4 conversations & OpenHermes (10K) & 5000 & 1.87 & Adapted \\
\bottomrule
\end{tabular}
\caption{Proximal NBGA fine-tuning of Qwen3.5-0.8B. All 24 blocks updated independently using $\delta_L$.}
\end{table}

\subsection{Compute Requirements}

\begin{itemize}
\item GPU: 24GB VRAM (RTX 4090)
\item Memory: $\sim$4 GB peak
\item Stage 1: $\sim$3 minutes (500 steps)
\item Stage 2: $\sim$3 minutes (500 steps)
\item Stage 3: $\sim$30 minutes (5000 steps)
\end{itemize}

\section{Scaling to 4B Parameters}

We demonstrate that NBGA scales to larger models by fine-tuning Qwen3.5-4B (32 blocks, 4B parameters) on OpenCodeInstruct (2K code generation examples). Unlike the 0.8B experiments, we use pure $\delta_L$ broadcast without the proximal term, enabled by a higher learning rate of $\eta = 10^{-2}$ that compensates for the gradient approximation error.

\begin{table}[h]
\centering
\begin{tabular}{lccc}
\toprule
Metric & Before NBGA & After NBGA & $\Delta$ \\
\midrule
Perplexity (val) & 3.20 & 2.16 & $-32\%$ \\
HumanEval pass@1 & 37.2\% & 42.7\% & $+5.5\%$ \\
\bottomrule
\end{tabular}
\caption{Qwen3.5-4B fine-tuned with NBGA on OpenCodeInstruct. 5000 steps, $\sim$30 minutes on RTX 4090.}
\end{table}

The HumanEval improvement is our first downstream task validation of NBGA at scale. The $\delta_L$ broadcast, while a coarse approximation, captures enough signal for the model to improve functional code generation. Training takes $\sim$30 minutes on a single RTX 4090 (5000 steps, 10 GB peak memory).

However, the HumanEval improvement peaks early (step 2900, 42.7\%) and degrades with continued training (step 5000, 32.3\%), despite perplexity remaining flat. The shared $\delta_L$ signal causes all layer weights to drift in correlated directions, overwriting useful structure when training continues past the optimal point.

\section{Controlled Comparison: TinyShakespeare}

We conduct controlled experiments on TinyShakespeare (character-level LM, 1M training tokens, $d=128$, conv+attn+MLP residual stack) to compare NBGA against full backprop, SPSA, and architectural variants under identical conditions.

\begin{table}[h]
\centering
\begin{tabular}{lcc}
\toprule
Method & Val PPL & $\Delta$ from backprop \\
\midrule
Full backprop & 6.50 & --- \\
NBGA + weight clamp ($\pm0.05$) & 6.75 & $+4\%$ \\
SPSA (per-layer) & 7.78 & $+20\%$ \\
NBGA + ReZero gates & 11.5 & $+77\%$ \\
NBGA + alternating SPSA & 11.75 & $+81\%$ \\
NBGA + layer diminishing ($1/N$) & 12.23 & $+88\%$ \\
NBGA + gated residual ($\sigma$) & 12.33 & $+90\%$ \\
Sparse predictive coding (local loss) & 18.38 & $+183\%$ \\
\bottomrule
\end{tabular}
\caption{TinyShakespeare controlled experiments. All methods use the same base architecture and training budget except where noted.}
\end{table}

Key findings from these ablations:
\begin{itemize}
\item \textbf{Weight clamping is essential.} Without it, weights grow, $\|J_f\|$ increases, and the $\delta_L \approx \delta_l$ approximation breaks down. Diminishing and gating cannot substitute for clamping.
\item \textbf{NBGA cannot learn ReZero-style gates.} When the residual path is initially closed ($\alpha=0$), the gradient $\delta_L \cdot f(h)$ vanishes because $f(h) \approx 0$ at initialization. The gate never opens.
\item \textbf{SPSA and NBGA are incompatible.} Alternating their updates produces conflicting gradient signals that hurt both methods.
\item \textbf{Local predictive losses trivialize.} Each layer learns to predict its input rather than learning task-relevant features, plateauing far above backprop performance.
\end{itemize}

\section{Limitations}

NBGA's approximation error grows with the spectral norm of the residual Jacobian $\|J_f\|$. For pre-trained models with large weights, the proximal term is essential; at 4B scale, even with high learning rate, the $\delta_L$ broadcast degrades downstream performance if training continues past a critical point (step $\sim$3000 for Qwen3.5-4B). The method is best used as a fast warm-start rather than a full training procedure.

NBGA cannot learn architectural gating mechanisms (ReZero, sigmoid gates, layer diminishing). When the residual path is initially attenuated, the gradient $\delta_L \cdot f(h)$ vanishes because $f(h) \approx 0$, preventing the gate from ever opening. Standard residual connections with weight clamping outperform all gated variants tested.

The method shares a deeper limitation with all chain-rule-free approaches: local predictive losses (as in predictive coding) trivialize to input reconstruction rather than learning task-relevant features, while global approximations (NBGA, SPSA) accumulate approximation error. No known method bridges the gap between local Hebbian plasticity and the expressiveness of backpropagation at scale.

\section{Conclusion}

NBGA demonstrates that residual networks can be fine-tuned without computing the chain rule between layers, using only the final-output gradient as a universal learning signal. A 752M-parameter, 24-block model was successfully instruction-tuned on Alpaca and OpenHermes data, with every block updated independently using $\delta_L$. Scaling to 4B parameters on OpenCodeInstruct yields a +5.5\% improvement on HumanEval pass@1, confirming that the $\delta_L$ broadcast carries sufficient signal for meaningful downstream improvement even at scale. The method is simple, parallelizable, memory-efficient, and runs on consumer hardware.

Controlled experiments on TinyShakespeare establish NBGA's position relative to backprop: within 4\% of full backprop perplexity when combined with weight clamping, outperforming SPSA and all architectural variants tested. However, the method degrades with extended training, cannot learn architectural gating, and fundamentally depends on small residual Jacobians. These findings map the boundaries of chain-rule-free training at scale and identify weight-constrained residual networks as the architecture best suited for gradient approximation methods.

\section*{Code}

All code and recipes available at: \url{https://github.com/Benzidrine/ngba_training}

\end{document}
